una vez alcanzada la superficie tenemos

\begin{equation*}
    s=e_2+me_1=0
\end{equation*}

entonces

\begin{equation*}
    e_2=-me_1
\end{equation*}

tenemos en cuenta que

\begin{equation*}
    \begin{aligned}
        \dot{e}_1 &= e_2 - f_m(x,t) = -me_1 - f_m(x,t)\\
    \end{aligned}
\end{equation*}

introducimos el cambio de variable

\begin{equation*}
    A_m=-m \quad \text{y} \quad B_m=-1
\end{equation*}

para tener

\begin{equation*}
    \dot{e}_1 = A_me_1 + B_mf_m(x,t)
\end{equation*}

se elige la función de Lyapunov

\begin{equation*}
    V_1 = e_1^Tee_1 \quad , \quad e>0
\end{equation*}

como $e_1$ es un escalar, tenemos que

\begin{equation*}
    V_1=pe_1^2
\end{equation*}

derivando

\begin{equation*}
    \begin{aligned}
        \dot{V}_1 &= \dot{e}_1^Tpe_1+\dot{e}_1^Tpe_1 \\
        &= e_1^T(PA_m+A_m^TP)e_1 + 2e_1^TPB_mf_m\\
    \end{aligned}
\end{equation*}

se busca garantizar que 

\begin{equation*}
    \dot{V}_1 \leq -\alpha V_1 + \gamma f_m^2
\end{equation*}

para $\alpha>0$ y $\gamma>0$. Usando la desiguladad de Young

\begin{equation*}
    2e_1^TPB_mf_m \leq \frac{-1}{\gamma}e_1^TPB_mB_m^TPe_1 + \gamma f_m^2
\end{equation*}

así

\begin{equation*}
    \begin{aligned}
        \dot{V}_1 &\leq e_1^T(PA_m+A_m^TP+PB_m\frac{1}{\gamma}B_m^TP)e_1 + \gamma f_m^2\\
        & \leq e_1^T(PA_m+A_m^TP+PB_m\frac{-1}{\gamma}B_m^TP+\alpha P)e_1 + \gamma f_m^2 \leq -\alpha P e_1^2
    \end{aligned}
\end{equation*}

lo que garantiza la epipsoide, entonce

\begin{equation*}
    \dot{V}_1 \leq e_1^T(A_mP+PA_m - PB_m\frac{-1}{\gamma}B_m^TP+\alpha P)e_1 + \gamma D_m^2
\end{equation*}

Usando el complemento de Schur, tenemos que

\begin{equation*}
    \begin{bmatrix}
        PA_m+A_m^TP+2\alpha P &  PB_m\\
        B_m^TP & -\gamma 
    \end{bmatrix} < 0 
\end{equation*}

permite que 

\begin{equation*}
    \dot{V}_1 \leq -\alpha V_1 + \gamma D_m^2
\end{equation*}
 \begin{equation*}
    V_1  \leq V_1(0) e^{-\alpha t} + \frac{\gamma D_m^2}{\alpha}(1-e^{-\alpha t})
 \end{equation*}


 de modo que 

\begin{equation*}
    \begin{aligned}
        \lim_{t \to \infty} V_1 &\leq \frac{\gamma D_m^2}{\alpha} \\
        \lim_{t \to \infty} e_1 &\leq \sqrt{\frac{\gamma D_m^2}{\alpha p}} \\
                                &\leq D_m \sqrt{\frac{\gamma}{\alpha P}}
    \end{aligned}
\end{equation*}

 Para el caso escalar la LMI es

 \begin{equation*}
    \begin{pmatrix}
        -(2m+\alpha)P & -P \\ -P & -\gamma
    \end{pmatrix} < 0
 \end{equation*}

 que equivale a 

 \begin{equation*}
    (-2m+\alpha)P+ \frac{P^2}{\gamma} <0
 \end{equation*}

como \(P>0\) y \(\gamma>0\)

 \begin{equation*}
    0<\alpha<2m
 \end{equation*}

 así 

 \begin{equation*}
    \gamma > \frac{P}{2m-\alpha}
 \end{equation*}

 para m fijo \begin{equation*}
    \alpha=m
 \end{equation*} y \begin{equation*}
    \gamma=\frac{P}{m}
 \end{equation*} , el radio  de la elipsoide es

 \begin{equation*}
    r=\frac{D_m}{m}
 \end{equation*}










 el infino del radio se obtiene en 

 \begin{equation*}
    r_{inf} \leq \frac{D_m}{m}
 \end{equation*}

 tambien es posible normalizar p, por ejemplo p=r implica

 \begin{equation*}
    \begin{bmatrix}
        -2m+\alpha & -1 \\
        -1 & -\gamma
    \end{bmatrix} < 0
 \end{equation*}

 que permite obtener una m variable a partir de un problema de factibilidad de la LMI

 En el caso n=2, el controlador es

 
 \begin{equation*}
    \begin{bmatrix}
        e_1&=r-y\\
        e_2&=\dot{r}-x_2\\
        s&=e_2+me_1\\
        u=me_2+KSign(s)
    \end{bmatrix}
 \end{equation*}

 donde \(K>R+D_u+mD_m\)

 Ejemplo: diseñe un controlador SMC-AEM para el sistema

 \begin{equation*}
    \begin{aligned}
        \dot{x}_1&=sin(x_1+x_2)+x_2\\
        \dot{x}_2&=cos(x_1+x_2)+u\\
        y&=x_1
    \end{aligned}
 \end{equation*}

